\section{Thiết lập thực nghiệm}
\label{sec:5-setup}

Mọi kết quả trong chương này được đo trên cùng một môi trường. Mỗi bảng kết quả ở các mục sau chỉ ghi các tham số khác với mục này.

\subsection{Phần cứng và phần mềm}
\label{sec:5-hardware}

{\small
\begin{longtblr}[
  caption = {Môi trường thực nghiệm},
  label = {tab:5-environment},
]{
  width = \linewidth,
  rowhead = 1,
  colspec = {X[1.4,l]X[3,l]}, rows = {abovesep=2pt, belowsep=2pt},
  hlines, vlines, row{1} = {font=\bfseries, bg=gray!20},
}
Hạng mục & Giá trị \\
CPU & Apple M5, 10 nhân (4 nhân hiệu năng, 6 nhân tiết kiệm), kiến trúc ARM64 \\
RAM của máy & 24~GB \\
Bộ nhớ cấp cho Docker & 11,7~GiB (OrbStack 2.2.5) \\
Ổ đĩa & SSD 926~GiB, còn trống khoảng 716~GiB \\
Hệ điều hành & macOS 27.2.1 \\
Docker Engine, Compose & Docker Engine 29.8.2, Compose 5.5.0 \\
Thành phần hệ thống & Theo Bảng~\ref{tab:4-versions} \\
Spark & \code{local[*]}, driver 2~GB, 4 nhân, 4 phân vùng shuffle \\
\end{longtblr}
}

\subsection{Dữ liệu thực nghiệm}
\label{sec:5-data}

Dữ liệu được sinh bởi bộ mô phỏng cảm biến (Mục~\ref{sec:4-application}) dựa trên dữ liệu danh mục \code{reference/reference\_data.json}: 4 trang trại, 8 thửa ruộng, 24 thiết bị, 4 loại cây trồng (lúa, ngô, cà chua, cà phê) và 5 chỉ số. Mỗi thiết bị báo cáo 5 giây một lần. Các giá trị ngẫu nhiên của bộ mô phỏng có hạt giống cố định, nhưng thời điểm đo lấy từ đồng hồ thật và mã sự kiện là UUID ngẫu nhiên, nên hai lần chạy cho chuỗi sự kiện tương tự chứ không giống hệt; vì vậy các kiểm tra ở Mục~\ref{sec:5-correctness} so với nhãn đi kèm từng sự kiện chứ không so với một tập dữ liệu cố định.

Bộ mô phỏng tiêm lỗi có chủ đích và đánh dấu mỗi sự kiện lỗi bằng trường \code{simulated\_fault}. Trường này đi theo bản ghi vào bronze (và nằm trong \code{raw\_payload} của bản ghi bị cách ly), nên có thể so số lỗi tiêm vào với số lỗi hệ thống phát hiện được mà không cần nhãn bên ngoài.

{\small
\begin{longtblr}[
  caption = {Cấu hình bộ mô phỏng (\texttt{.env.example})},
  label = {tab:5-simulator},
]{
  width = \linewidth,
  rowhead = 1,
  colspec = {X[1.8,l]X[1.2,l]X[2.6,l]}, rows = {abovesep=2pt, belowsep=2pt},
  hlines, vlines, row{1} = {font=\bfseries, bg=gray!20},
}
Tham số & Mặc định & Ý nghĩa \\
\code{SIM\_SEED} & 42 & Hạt giống ngẫu nhiên \\
\code{SIM\_RATE\_HZ} & 1 & Số nhịp sinh dữ liệu mỗi giây \\
\code{SIM\_FAULT\_RATE} & 0,05 & Xác suất ở mỗi nhịp có một thiết bị ngẫu nhiên được lên lịch tiêm lỗi \\
\code{SIM\_FAULTS} & 7 loại & \code{null\_value}, \code{out\_of\_range}, \code{implausible\_jump}, \code{duplicate}, \code{unknown\_device}, \code{unexpected\_metric}, \code{malformed} \\
\code{SIM\_GAP\_RATE} & 0,002 & Xác suất ở mỗi nhịp có một thiết bị ngẫu nhiên bắt đầu im lặng \\
\code{SIM\_GAP\_SECONDS\_MIN/MAX} & 45 / 90 & Độ dài khoảng im lặng (giây) \\
\code{SIM\_DURATION\_SECONDS} & 0 & Thời gian chạy một lần thực nghiệm (0 = chạy liên tục; mỗi benchmark tự bật, tắt bộ mô phỏng và ghi độ dài cửa sổ đo vào CSV) \\
\end{tblr}
\end{table}

\subsection{Tham số của đường xử lý}
\label{sec:5-params}

{\small
\begin{longtblr}[
  caption = {Tham số mặc định của tác vụ streaming và cảnh báo},
  label = {tab:5-params},
]{
  width = \linewidth,
  rowhead = 1,
  colspec = {X[2,l]X[1,l]X[2.6,l]}, rows = {abovesep=2pt, belowsep=2pt},
  hlines, vlines, row{1} = {font=\bfseries, bg=gray!20},
}
Tham số & Giá trị & Ý nghĩa \\
\code{TRIGGER\_INTERVAL} & 30 s & Chu kỳ micro-batch (benchmark \#1 đổi thành 10 s và 60 s) \\
\code{MAX\_OFFSETS\_PER\_TRIGGER} & 20\,000 & Số sự kiện tối đa mỗi lô \\
\code{WATERMARK\_S} & 3\,600 & Cửa sổ chấp nhận dữ liệu đến muộn \\
\code{LATE\_AFTER\_S} & 600 & Ngưỡng gắn cờ \code{is\_late} \\
\code{ALERT\_DEBOUNCE\_BATCHES} & 2 & Số lô vi phạm liên tiếp trước khi cảnh báo \\
\code{ALERT\_COOLDOWN\_S} & 300 & Khoảng tối thiểu giữa hai cảnh báo cùng điều kiện \\
\code{DEVICE\_OFFLINE\_MULTIPLIER} & 4 & Im lặng quá 4 lần chu kỳ báo cáo (20 s) thì coi là mất kết nối \\
\code{BATTERY\_LOW\_PCT} & 20 & Ngưỡng pin yếu (\%) \\
\code{SIGNAL\_POOR\_MAX} & 0,3 & Ngưỡng sóng yếu \\
\end{longtblr}
}

\subsection{Quy trình chạy lại}
\label{sec:5-procedure}

Các benchmark chạy trên tập dữ liệu tích lũy của hệ thống (khoảng 41 nghìn số đo trong bronze trên 12 ngày có dữ liệu, cộng dữ liệu sinh ra trong lúc đo), không xóa sạch giữa các lần đo, vì một số phép đo (gold, time travel, compaction) cần lịch sử. Benchmark \#2--\#4 tự tạo dữ liệu tổng hợp của mình trong schema \code{bench}. Trình tự đã dùng:

\begin{lstlisting}[language=bash, caption={Trình tự chạy các benchmark và kiểm tra}, label={lst:5-procedure}]
just up && just up-app        # all profiles, data volumes kept
just bench-lookups            # #2  (stream and simulator stopped)
just bench-storage            # #3
just bench-compaction         # #4
just bench-trigger            # #1  restarts spark-stream per interval
just bench-replay             # #5  kills/stops spark-stream, checks Kafka vs bronze
just stop-stream && just batch-run   # dbt + tests + compaction (bronze, ops)
just check-correctness --since <time of the simulator fix>   # section 5.2
just run-stream-bg && just bench-web   # #6 (stream and simulator running)
\end{lstlisting}

Kết quả của từng benchmark được lưu trong \code{benchmarks/results/*.csv}, với cấu hình của lần đo (máy, phiên bản, tải, cửa sổ đo, trigger) ghi ở các dòng đầu tệp.
